\chapter{Verification conventions and tolerances}\label{ch:conv}

\section{A template for a convergence study}\label{sec:conv-template}

A convergence study answers one question: does the error of a discretisation fall at the rate that the method promises? The studies of this book follow one template, and \Cref{ch:elem} to \Cref{ch:cms} use it where a mesh, a time step or a basis size is varied (the package provides \anchor{fea_engine.convergence.run_study} and \anchor{fea_engine.convergence.observed_order}).

\begin{enumerate}
\item \textbf{Fix the reference.} An exact solution, a closed form, or an independent computation (a different solver or a much finer model). A reference produced by the method under test shows only self-consistency.
\item \textbf{Fix the quantity and the norm.} A global norm (relative $L_2$ error of the displacement) and a local quantity (a tip deflection, a reaction, a peak stress); the definition of the error is stated with the result.
\item \textbf{Vary one parameter} over at least three values in the asymptotic range (mesh size $h$, step $\Delta t$, basis size $r$), keeping everything else fixed.
\item \textbf{Compute the observed order} $p=\log(e_1/e_2)/\log(h_1/h_2)$ from successive pairs and compare it with the theoretical order; a rate that drifts, or an error that stops falling, is reported and explained (rounding floor, reference error, a floor set by the basis).
\item \textbf{State what was held fixed and what was not tested}: the load, the boundary conditions, the other discretisation parameters, the range of the parameter.
\end{enumerate}

For reduced-order models the parameter of step 3 is the basis size, there is no asymptotic order to compare, and the template becomes: report the error on \emph{held-out} parameters and not on the training set, and report both (\Cref{ch:pod,ch:contact}).

\section{Tolerances used in the book}\label{sec:conv-tol}

\Cref{tab:conv-tol} lists the numerical tolerances that decide whether a solve is called converged, as they appear in the listings of this book and in the defaults of the code that Chapters~\ref{ch:nonint} and~\ref{ch:hyper} call. A tolerance is a setting, not a measured accuracy: a result converged to $10^{-10}$ in the residual is accurate to the discretisation, not to $10^{-10}$.

\begin{table}[ht]
\centering\small
\begin{tabularx}{\linewidth}{@{}>{\raggedright\arraybackslash}p{0.38\linewidth} >{\raggedright\arraybackslash}p{0.2\linewidth} X@{}}
\toprule
Where & Value & Meaning \\
\midrule
nonlinear static drivers (listings of Chapters~\ref{ch:nonlin}, \ref{ch:batch}) & \code{tol} $10^{-9}$ to $10^{-12}$ & relative residual at which Newton stops; \code{max_iter} 40--60 \\
\code{lcp_active_set} & $10^{-12}$ (default) & tolerance of the active-set linear complementarity solver \\
\code{ContactROM.solve_dynamic} & $10^{-4}$ (relative) & a reduced step whose contact conditions are violated by more than this is reported as a failure; a guard, not an accuracy claim \\
\code{solve_reduced_contact_newton} & \code{tol} $10^{-9}$, \code{max_iter} 60 & residual and step of the reduced contact Newton iteration \\
\code{ecsw_weights} & \code{tol} $10^{-6}$ (default) & relative training residual at which element selection stops (Chapter~\ref{ch:hyper} sweeps $10^{-2}$ to $10^{-8}$) \\
\code{nonneg_dual_basis} & \code{tol} $10^{-9}$, 500 iterations & non-negative factorisation stop \\
\code{rank_deficiency} & $10^{-10}$ (relative) & singular values below this fraction of the largest count as zero \\
success criteria of the experiments & 1\,\% (Chapter~\ref{ch:pod}), 2\,\% (Chapter~\ref{ch:nonint}) & held-out error and frequency error used to call a basis or an identification adequate; chosen by the author, not by a standard \\
\bottomrule
\end{tabularx}
\caption{Tolerances and thresholds.}\label{tab:conv-tol}
\end{table}

\section{What a reported ROM number states}\label{sec:conv-rom}

A speed-up or error figure of a reduced-order model is meaningful only with its conditions. Each such number in this book should state:
\begin{itemize}
\item the full-order model, its size, its solver and whether the baseline is dense or sparse;
\item the \emph{offline} cost (snapshots, basis, training, preprocessing), the \emph{online} cost per query, and, if the answer has to be lifted to the full space, the \emph{reconstruction} cost, each separately;
\item whether a quoted speed-up includes the preprocessing: a speed-up of the online solve alone is not a speed-up of the workflow, and the break-even number of queries (offline cost divided by saving per query) is what connects them;
\item the error on held-out parameters, with its definition, next to a \emph{local} quantity of engineering interest (a reaction, a peak stress or displacement), because a small global error can hide a large local one.
\end{itemize}
Chapters~\ref{ch:pod} to~\ref{ch:hyper} report online speed-ups against dense baselines on small models, separate offline and online costs (Chapter~\ref{ch:pod} adds break-even counts), and report global errors on held-out parameters. They do \emph{not} report reconstruction time separately, and local quantities are reported only where a chapter's quantity of interest is itself local (for instance the contact force of Chapter~\ref{ch:hyper}); this is a limitation of the book, listed in the plan.
